\label{der:2-2}
若$A$可逆，$A\vx=0$左乘$A\inv$给$\vx=0$。反过来，对$A$作初等行变换得到阶梯形$U=EA$，其中$E$是可逆初等矩阵的乘积。每种行操作都有逆操作：换行再换回，倍乘$c\ne0$的逆是倍乘$\frac{1}{c}$，倍加$c$的逆是倍加$-c$，所以$A\vx=0$与$U\vx=0$同解。

若方阵的主元少于$n$，至少有一个自由变量；令一个自由变量为1、其余为0，回代得到非零齐次解，与假设矛盾。因此每列均为主元列，每行也非零，继续把主元归一并消去上方元素，存在初等矩阵乘积$F$使$FA=I$。于是$A=F\inv$，其逆为$F$。这同时证明行等价于$I$、可逆、满秩、无非零齐次解彼此等价，也证明$A$能写成有限个初等矩阵的乘积。反向，初等矩阵均可逆，其有限乘积可逆。

由乘积行列式，$A$可逆推出$\det A\det A\inv=1$，所以行列式非零；若行列式非零，伴随恒等式给
\[A\dfrac{\adj A}{\det A}=\dfrac{\adj A}{\det A}A=I,\]
故可逆。若两个逆为$B,C$，则$B=B(AC)=(BA)C=C$，逆唯一。

对增广矩阵，每个左乘初等矩阵同时作用于两块：$[A\mid I]\mapsto[FA\mid F]$。当左块为$I$时$F=A\inv$，所以右块恰为逆；这解释了同步行变换法而不只是给出算法。
